\section{Conclusions and Future Work}
\label{sec:conclusions}

This report addressed change impact analysis for Terraform: given an edited block, which other
blocks will change with it, and which blocks will be added? We mined a dataset from 53 open-source
repositories and learned, with a self-implemented Genetic Algorithm, how to combine twelve signals
from the repository's history, reference graph and text. \textbf{RQ1}: the Genetic Algorithm is a
feasible and reliable learner for this task: it learns from the labels, converges within about 60
generations and gives almost identical results across runs and settings, although random search
with the same budget comes close. \textbf{RQ2}: the learned ranking clearly outperforms
history-based rule mining (ROSE) and equal weights, mainly by placing new block types and by
following references from removed blocks, and the three signal families are complementary. The
\tool{} tool makes the learned model usable on any local Terraform repository.

Future work follows from the limits above. Non-linear rankers, such as gradient-boosted ranking
models or Genetic Programming (GP) that evolves the scoring formula, could capture the interactions a
linear score misses and give the evolutionary search a harder problem; a simpler step is separate
weights per seed action. An independent sample of new repositories would test the design choices
as well as the weights, and other cloud providers and IaC languages would test generality. For
everyday use, caching the mined history would make repeated recommendations instant.
